% 12-resultats-quantitatifs.tex — Résultats 2/4 : résultats quantitatifs (≈ 7 pages).

\chapter{Résultats quantitatifs}
\label{chap:resultats-quantitatifs}

\section{Participation et données recueillies}

\begin{arediger}{1 page}
Nombre de participants à chaque étape, taux de réponse, données exclues et
pourquoi. Tableau descriptif.
\end{arediger}

\section{Compréhension conceptuelle}

\begin{arediger}{3 pages}
Scores pré-test et post-test (histogrammes ou boîtes à moustaches en
pgfplots), gain normalisé global et par thème (énergie, mouvement circulaire,
réaction normale), comparaison avec le groupe témoin, tests et tailles
d'effet. Analyse item par item des questions sur le looping et le seau.
\end{arediger}

\section{Utilisation des simulations et du mode jeu}

\begin{arediger}{2 pages}
Temps d'utilisation, exercices les plus utilisés, nombre d'essais et de
réussites par défi, paliers atteints, variable libre la plus difficile
(\(v_0\), \(L\) ou \(m\)). Croiser avec les gains.
\end{arediger}

\section{Perception des étudiants}

\begin{arediger}{1 page}
Résultats du questionnaire de Likert : graphique en barres empilées par item.
\end{arediger}
